
% ============================================================
% Research-track patch: Taylor-column finite-thickness and parity-delay benchmark
% ============================================================
\subsection{Research Track: Taylor-Column Analogues for Finite-Thickness Swirl Strings}
\label{sec:taylor_column_finite_thickness_swirl_strings}

\providecommand{\rhoF}{\rho_{\!f}}
\providecommand{\vSwirl}{\mathbf{v}_{\!\boldsymbol{\circlearrowleft}}}
\providecommand{\rc}{r_c}
\providecommand{\GammaK}{\Gamma_K}
\providecommand{\GammaZero}{\Gamma_0}
\providecommand{\Om}{\Omega}
\providecommand{\Ht}{\mathcal{H}_{\times}}

\paragraph{Status.}
This subsection is \textbf{Research Track}.  The Taylor--Proudman,
Hill-vortex, relative-helicity, inertial-wave, and Feynman--Onsager statements
are orthodox rotating-fluid or superfluid mechanics
\cite{Proudman1916,Taylor1922,Batchelor1967,Greenspan1968,Hill1894,Onsager1949,Feynman1955}.
Their SST use is restricted to analogue benchmarking, scale-separation
diagnostics, and falsification tests.  No macroscopic rotating-tank parameter is
identified with the SST canonical core scale.

\subsubsection{Background kinematics and resolved scale separation}

Consider a rotating cylindrical background of height
\(H=1\,\mathrm m\), radius \(R=0.25\,\mathrm m\), and constant angular velocity
\[
\boldsymbol{\Omega}=\Omega\,\hat{\mathbf z},
\qquad
\Omega=1\,\mathrm{s^{-1}}.
\]
The unperturbed velocity and vorticity fields are
\[
\mathbf v_b=\Omega r\,\hat{\boldsymbol\theta},
\qquad
\boldsymbol\omega_b=\nabla\times\mathbf v_b
=2\Omega\,\hat{\mathbf z}.
\]
The background has zero local kinetic helicity density:
\[
\mathbf v_b\cdot\boldsymbol\omega_b
=(\Omega r\,\hat{\boldsymbol\theta})\cdot(2\Omega\,\hat{\mathbf z})=0.
\]

Let \(a\) denote the resolved separatrix radius of a macroscopic vortex
structure, such as a vortex ring, Hill-type vortex, or trefoil vortex tube.
Let \(r_c\) denote the SST canonical core radius.  The resolved-scale guardrail is
\[
\boxed{a\neq r_c.}
\]
Using
\[
r_c=1.40897017\times10^{-15}\,\mathrm m,
\qquad
\vSwirl=1.09384563\times10^6\,\mathrm{m\,s^{-1}},
\]
the canonical core-scale circulation is
\[
\Gamma_0=2\pi r_c\vSwirl
=9.68361920\times10^{-9}\,\mathrm{m^2\,s^{-1}}.
\]

For a macroscopic separatrix radius
\[
a=\frac{R}{4}=0.0625\,\mathrm m,
\]
the background circulation around the same radius is
\[
\Gamma_b(a)=\oint_{r=a}\mathbf v_b\cdot d\boldsymbol\ell
=2\pi\Omega a^2.
\]
This defines the dimensionless synchronization parameter
\[
\boxed{
\chi_\Omega=\frac{\Gamma_K}{2\pi\Omega a^2}.
}
\]
For \(\Gamma_K=\Gamma_0\) and \(\Omega=1\,\mathrm{s^{-1}}\),
\[
2\pi\Omega a^2=2.45436926\times10^{-2}\,\mathrm{m^2\,s^{-1}},
\]
so that
\[
\boxed{
\chi_\Omega=3.94546141\times10^{-7}.
}
\]
Thus a macroscopic separatrix equipped only with the canonical core circulation
is kinematically far below synchronization with a \(1\,\mathrm{s^{-1}}\)
tank-scale rotation.

\paragraph{Dimensional check.}
\[
[\Gamma_K]=\mathrm{m^2\,s^{-1}},
\qquad
[\Omega a^2]=\mathrm{s^{-1}m^2}=\mathrm{m^2\,s^{-1}},
\]
so \(\chi_\Omega\) is dimensionless.

\subsubsection{Separatrix response and relative helicity}

For a slowly translating finite-thickness structure with axial speed \(U\), the
Rossby number is
\[
Ro=\frac{U}{2\Omega a}.
\]
In the low-Rossby regime \(Ro\ll1\), the leading-order Taylor--Proudman
constraint gives
\[
\frac{\partial\mathbf u}{\partial z}\approx0.
\]
Hence, to leading order, a no-through-flow separatrix behaves externally as a
columnar obstruction.  This does not imply that the internal topology is
sphere-like; it only states that the exterior flow sees a columnar separatrix
constraint.

For an incompressible closed cross-section, the axial flux balance is
\[
\boxed{
\int_0^R2\pi r\,u_z(r)\,dr=0.
}
\]
A simple two-zone model is
\[
u_z(r)=
\begin{cases}
U, & 0\le r<a,\\[1mm]
-U\dfrac{a^2}{R^2-a^2}, & a<r\le R.
\end{cases}
\]
Indeed,
\[
\int_0^a2\pi rU\,dr
+
\int_a^R2\pi r\left(-U\frac{a^2}{R^2-a^2}\right)dr
=\pi a^2U-\pi a^2U=0.
\]

For a vortex ring with circulation \(\Gamma_K\) and ring surface \(S_K\), a
useful diagnostic is the relative or mutual helicity with the rotating
background \cite{Moffatt1969}:
\[
\boxed{
\mathcal H_{\times}
=2\Gamma_K\int_{S_K}\boldsymbol\omega_b\cdot d\mathbf S.
}
\]
Because
\[
\boldsymbol\omega_b=2\Omega\hat{\mathbf z},
\qquad
\int_{S_K}\hat{\mathbf z}\cdot d\mathbf S=\pi a^2,
\]
one obtains
\[
\boxed{
\mathcal H_{\times}=4\pi\Omega a^2\Gamma_K.
}
\]
For \(\Gamma_K=\Gamma_0\),
\[
\boxed{
\mathcal H_{\times,0}
=4.75343546\times10^{-10}\,\mathrm{m^4\,s^{-2}}.
}
\]

\paragraph{Dimensional check.}
\[
[\mathcal H_{\times}]
=[\Omega a^2\Gamma_K]
=\mathrm{s^{-1}m^2\,m^2s^{-1}}
=\mathrm{m^4\,s^{-2}}
=[\Gamma_K]^2.
\]

\subsubsection{Hill-vortex sphere analogue}

A vortex-ring dipole with a closed separatrix may be approximated, in the
continuous Euler limit, by Hill's spherical vortex \cite{Hill1894}.  In the
frame where the undisturbed fluid at infinity is at rest, the kinetic energy
splits into an external irrotational contribution and an internal rotational
contribution:
\[
E_{\rm out}=\frac{1}{3}\pi\rho_{\!f}a^3U^2,
\qquad
E_{\rm in}=\frac{23}{21}\pi\rho_{\!f}a^3U^2.
\]
Therefore
\[
\boxed{
E_{\rm tot}=E_{\rm out}+E_{\rm in}
=\frac{10}{7}\pi\rho_{\!f}a^3U^2.
}
\]
This result is used only as a continuum benchmark for finite-thickness
separatrix energetics.  It does not identify a trefoil vortex with a
featureless sphere; the equivalence is external and leading-order only.

\paragraph{Numerical reference for \(U=1\,\mathrm{m\,s^{-1}}\).}
With
\[
\rho_{\!f}=7.0\times10^{-7}\,\mathrm{kg\,m^{-3}},
\qquad
a=0.0625\,\mathrm m,
\]
one obtains
\[
E_{\rm out}=1.78964425\times10^{-10}\,\mathrm J,
\]
\[
E_{\rm in}=5.88025969\times10^{-10}\,\mathrm J,
\]
and
\[
E_{\rm tot}=7.66990394\times10^{-10}\,\mathrm J.
\]

\subsubsection{Axial force density from vorticity-pressure gradients}

If the local background rotation varies along the axis, \(\Omega=\Omega(z)\),
ordinary Euler/Bernoulli balance gives
\[
\frac{1}{\rho_{\!f}}\frac{\partial p}{\partial r}
=\Omega(z)^2r.
\]
Integrating from the axis to the wall yields
\[
p(0,z)=p(R,z)-\frac{1}{2}\rho_{\!f}\Omega(z)^2R^2.
\]
If \(p(R,z)\) is fixed or externally controlled, the axial force density on the
central region is
\[
\boxed{
f_z=-\frac{\partial p(0,z)}{\partial z}
=\frac{1}{2}\rho_{\!f}R^2
\frac{\partial}{\partial z}\left[\Omega(z)^2\right].
}
\]
This force is part of the ordinary vorticity-pressure impulse budget:
\[
\boxed{
\mathbf F_z\subset\mathbf F_{\omega}.
}
\]
It must be analytically subtracted before any residual SST-specific force claim
is considered.

\paragraph{Dimensional check.}
\[
[\rho_{\!f}R^2\partial_z\Omega^2]
=\mathrm{kg\,m^{-3}}\,\mathrm{m^2}\,\mathrm{m^{-1}s^{-2}}
=\mathrm{kg\,m^{-2}s^{-2}}
=\mathrm{N\,m^{-3}}.
\]
Thus \(f_z\) has the unit of force density.

\subsubsection{Trefoil dynamics and He-II falsification}

A macroscopic trefoil vortex is not a sphere internally.  Its internal geometry
carries writhe, twist, finite-thickness contact structure, and self-induced
Biot--Savart dynamics.  In water-like classical fluids, hydrofoil-generated
trefoil vortices are expected to deform strongly, redistribute helicity, and
undergo reconnection \cite{KlecknerIrvine2013,KoplikLevine1993}.

In orthodox rotating superfluid helium, the background rotation is carried by a
Feynman--Onsager lattice of quantized vortex lines \cite{Onsager1949,Feynman1955}.
For \({}^4\mathrm{He}\),
\[
\kappa_4=\frac{h}{m_4}=9.9693\times10^{-8}\,\mathrm{m^2\,s^{-1}}.
\]
The vortex-line areal density is
\[
n_v=\frac{2\Omega}{\kappa_4}.
\]
For \(\Omega=1\,\mathrm{s^{-1}}\),
\[
\boxed{
n_v=2.0062\times10^7\,\mathrm{m^{-2}}.
}
\]
The number of quantized vortex lines crossing the resolved separatrix area
\(\pi a^2\) is
\[
N_v=n_v\pi a^2.
\]
For \(a=0.0625\,\mathrm m\),
\[
\boxed{
N_v=2.4619\times10^5.
}
\]

The orthodox He-II expectation is therefore a high reconnection rate,
Kelvin-wave cascade, partial helicity redistribution, and likely unknotting or
decay of the trefoil.  By contrast, the ideal SST continuum limit assumes
inviscid, incompressible flow with no reconnecting cuts, so topology is preserved
unless an explicit reconnection or defect-bridge sector is added.  The
disagreement defines a falsification domain:
\[
\boxed{\text{He-II quantized vortex lattice: reconnection-enabled decay,}}
\]
\[
\boxed{\text{ideal SST continuum: topology-preserving transport.}}
\]

\subsection{Research Track: Rotating-Fluid Delay and Swirl-Clock Parity}
\label{sec:rotating_fluid_swirl_clock_parity}

\paragraph{Status.}
This subsection is \textbf{Research Track}.  It uses orthodox inertial-wave and
rotating-vorticity relations as analogue benchmarks and then records the
SST-compatible parity split between odd tracer observables and even clock-rate
observables.  No dimensionally invalid fine-structure normalization is used.

\subsubsection{Arrival delay of the surface signal: inertial-wave delay}

For a small perturbation in a uniformly rotating inviscid incompressible fluid,
the inertial-wave dispersion relation is
\[
\omega=2\Omega\frac{k_z}{k},
\qquad
k^2=k_r^2+k_z^2.
\]
The vertical group velocity is
\[
c_{g,z}=\frac{\partial\omega}{\partial k_z}
=2\Omega\frac{k_r^2}{k^3}.
\]
Therefore the arrival time of a disturbance over height \(H\) is
\[
\boxed{
t_{\mathrm{arr}}
=\frac{H}{c_{g,z}}
=\frac{Hk^3}{2\Omega k_r^2}.
}
\]
Thus
\[
\boxed{t_{\mathrm{arr}}\propto\Omega^{-1}.}
\]
This is a propagation delay of the rotating-fluid disturbance.  It is not a
swirl-clock time-dilation effect.

\paragraph{Dimensional check.}
\[
\left[\frac{Hk^3}{2\Omega k_r^2}\right]
=\frac{\mathrm{m}\,\mathrm{m^{-3}}}{\mathrm{s^{-1}m^{-2}}}
=\mathrm{s}.
\]

\subsubsection{Clock delay from local swirl speed: even-in-direction effect}

In the SST swirl-clock interpretation, the local proper-time rate associated
with a resolved azimuthal swirl speed \(u_\theta\) is taken in the dimensionally
consistent form
\[
\boxed{
\frac{d\tau}{dt}
=\sqrt{1-\frac{u_\theta^2}{c^2}}.
}
\]
For \(|u_\theta|\ll c\),
\[
\frac{d\tau}{dt}
=1-\frac{1}{2}\frac{u_\theta^2}{c^2}
+\mathcal O\!\left(\frac{u_\theta^4}{c^4}\right).
\]
The accumulated proper-time deficit over a laboratory interval \(T\) is
therefore
\[
\Delta\tau-T
=\int_0^T\left(\sqrt{1-\frac{u_\theta(t)^2}{c^2}}-1\right)dt
\approx
-\frac{1}{2c^2}\int_0^Tu_\theta(t)^2\,dt.
\]
Because the dependence is quadratic,
\[u_\theta\rightarrow-u_\theta
\quad\Longrightarrow\quad
u_\theta^2\rightarrow u_\theta^2,
\]
so
\[
\boxed{
\Delta\tau(+u_\theta)=\Delta\tau(-u_\theta)<T.
}
\]
Both swirl directions therefore produce clock slowing.  Directional reversal
changes the sign of tracer rotation, but not the sign of the clock-rate deficit.

\subsubsection{Vorticity response to vertical displacement}

In the low-Rossby rotating-fluid limit, the linearized vorticity equation gives
\[
\partial_t\boldsymbol\omega
\approx
2(\boldsymbol\Omega\cdot\nabla)\mathbf u.
\]
For an axisymmetric vertical displacement field
\[
\mathbf u=w(r,z,t)\,\hat{\mathbf z},
\qquad
w=\partial_t\xi,
\]
the vertical vorticity response is
\[
\partial_t\omega_z
=2\Omega\,\partial_z w
=2\Omega\,\partial_z\partial_t\xi.
\]
Integrating in time with zero initial perturbation vorticity gives
\[
\boxed{
\omega_z(r,z,t)=2\Omega\,\partial_z\xi(r,z,t).
}
\]
Thus column stretching and compression generate opposite-signed vertical
vorticity.

For a Gaussian displacement profile,
\[
\xi(r,z,t)
=Z(t)\exp\left[-\frac{r^2+z^2}{a^2}\right],
\]
one obtains
\[
\omega_z(r,z,t)
=-\frac{4\Omega Z(t)z}{a^2}
\exp\left[-\frac{r^2+z^2}{a^2}\right].
\]
Hence the induced vorticity is odd in \(z\):
\[
\boxed{
\omega_z(r,-z,t)=-\omega_z(r,z,t).
}
\]

\subsubsection{Resolved swirl velocity from induced vorticity}

For an axisymmetric azimuthal velocity \(u_\theta(r,z,t)\),
\[
\omega_z=\frac{1}{r}\frac{\partial}{\partial r}\left(r u_\theta\right).
\]
Using the Gaussian vorticity source above gives
\[
\boxed{
u_\theta(r,z,t)
=-\frac{2\Omega Z(t)z}{a^2}
e^{-z^2/a^2}
\frac{1-e^{-r^2/a^2}}{r}.
}
\]
Therefore
\[u_\theta(r,-z,t)=-u_\theta(r,z,t),
\qquad
u_\theta(r,-z,t)^2=u_\theta(r,z,t)^2.
\]
This establishes the parity split:
\[
\boxed{
\text{tracer angle is odd in swirl direction,}
\qquad
\text{clock deficit is even in swirl direction.}
}
\]

\subsubsection{C1--C4 conclusions}

\paragraph{C1. Inertial-wave benchmark.}
The delayed push--pull surface response in a rotating tank is explained, to
leading order, as an inertial-wave arrival delay:
\[
\boxed{
t_{\mathrm{arr}}=\frac{Hk^3}{2\Omega k_r^2}.
}
\]
It scales as \(t_{\mathrm{arr}}\propto\Omega^{-1}\), and therefore disappears
as a rotating-fluid effect when the background rotation is removed.

\paragraph{C2. Vorticity generation by vertical forcing.}
A vertical displacement in a rotating incompressible fluid produces vertical
vorticity according to
\[
\boxed{
\omega_z=2\Omega\,\partial_z\xi.
}
\]
Stretching and compression of the rotating column generate opposite signs of
relative vorticity above and below the symmetry plane.

\paragraph{C3. Odd tracer response and cancellation over symmetric cycles.}
The induced azimuthal velocity satisfies
\[u_\theta(r,-z,t)=-u_\theta(r,z,t).
\]
Hence tracer angles reverse under reflection or reversed forcing.  For an ideal
symmetric up--down cycle, the leading-order angular displacement cancels:
\[
\boxed{
\Delta\theta_{\mathrm{cycle}}\simeq0
\qquad
\text{to leading order in }Ro.
}
\]

\paragraph{C4. Even clock-rate response and accumulated delay.}
The swirl-clock deficit depends on \(u_\theta^2\), not on \(u_\theta\).  Therefore
\[
\boxed{
\Delta\tau(+u_\theta)=\Delta\tau(-u_\theta)<T.
}
\]
Equivalently,
\[
\boxed{
\Delta\tau-T
\approx
-\frac{1}{2c^2}\int_0^Tu_\theta(t)^2\,dt.
}
\]
Thus a symmetric cycle can have zero net tracer rotation while still producing a
nonzero accumulated clock deficit.

\paragraph{Core-scale guardrail.}
The laboratory rotation rate \(\Omega\) must not be identified with the SST
canonical core turnover rate.  The latter is
\[
\Omega_0=\frac{\vSwirl}{r_c}.
\]
Using the canonical constants gives
\[
\boxed{
\Omega_0=7.76344066\times10^{20}\,\mathrm{s^{-1}}.
}
\]
Therefore
\[
\boxed{
\Omega_{\mathrm{lab}}\neq\Omega_0.
}
\]
The rotating-tank experiment is a resolved analogue benchmark, not a direct
identification of macroscopic rotation with the canonical SST core scale.

\paragraph{Summary.}
The experiment separates two parities:
\[
\boxed{
\theta(+u_\theta)=-\theta(-u_\theta),
}
\]
but
\[
\boxed{
\Delta\tau(+u_\theta)=\Delta\tau(-u_\theta).
}
\]
This provides a falsifiable benchmark: angular tracer motion is direction
sensitive, whereas the clock-rate or energy-density deficit is direction
independent at leading quadratic order.

\paragraph{Child-level analogy.}
A knotted vortex in a rotating fluid is like a small spinning propeller inside a
slowly rotating elevator shaft.  From outside, it may push fluid like a round
object moving up and down, but inside the object the strings are tied into a
knot.  Turning the knot the other way reverses the visible swirl, but the amount
of busy motion stored in the fluid still slows the clock in both directions.


\bibitem{Taylor1922}
G.~I. Taylor,
``The motion of a sphere in a rotating liquid,''
\emph{Proceedings of the Royal Society of London A} \textbf{102} (1922), 180--189.
DOI: \href{https://doi.org/10.1098/rspa.1922.0079}{10.1098/rspa.1922.0079}.

\bibitem{Proudman1916}
J.~Proudman,
``On the motion of solids in a liquid possessing vorticity,''
\emph{Proceedings of the Royal Society of London A} \textbf{92} (1916), 408--424.
DOI: \href{https://doi.org/10.1098/rspa.1916.0026}{10.1098/rspa.1916.0026}.

\bibitem{Hill1894}
M.~J.~M. Hill,
``On a spherical vortex,''
\emph{Philosophical Transactions of the Royal Society of London A} \textbf{185} (1894), 213--245.
DOI: \href{https://doi.org/10.1098/rsta.1894.0006}{10.1098/rsta.1894.0006}.

\bibitem{Greenspan1968}
H.~P. Greenspan,
\emph{The Theory of Rotating Fluids},
Cambridge University Press (1968).
ISBN: 978-0521051477.

